
% ================================================================
% CHAPTER 3: REQUIREMENT ENGINEERING
% ================================================================

\section{Feasibility Study}

\subsection{Technical Feasibility}

The system is built entirely on mature, open-source libraries that have been
in active use for years. Neo4j Community Edition handles graph storage, Ollama
runs the LLM locally, and the sentence-transformers library handles embeddings.
All of these install via \texttt{pip} and work on standard hardware, a laptop
with 16\,GB RAM is sufficient. A dedicated GPU would speed things up but is not
necessary to run the system.

\subsection{Operational and Economic Feasibility}

The pipeline runs through a single command:
\begin{verbatim}
python enhanced_main.py --corpus my_documents.json
\end{verbatim}
No configuration files need to be written, and no cloud account is needed.
Every component used, Neo4j, Ollama, Qwen2.5, and all Python packages, is
free and open-source. There are no per-query costs.

\subsection{Legal Feasibility}

All components are published under OSI-approved licences (GPL v3, MIT, Apache
2.0) that permit academic use. Since no data leaves the local machine, the
system is compatible with institutional data-handling policies.

% ----------------------------------------------------------------
\section{Requirement Elicitation and Analysis}

\subsection{Objective}

The primary objective of the requirements phase was to formally capture what the
system must do (functional requirements) and what quality properties it must
exhibit (non-functional requirements), independently of implementation decisions.

\subsection{Methods Used for Elicitation}

Requirements were elicited through three methods:

\begin{enumerate}
    \item \textbf{Paper analysis}: A detailed reading of arXiv:2511.10375 was
    used to extract the algorithm specifications of Modules A, B, and C as
    implicit functional requirements. If the paper says a module does X, the
    implementation must be testably verified to do X.
    \item \textbf{Gap analysis}: The v4 output was compared against known
    failure modes of standard RAG (temporal blindness, no confidence score,
    no audit trail) to identify the extensions needed for v5.
    \item \textbf{Prototype evaluation}: An early prototype was run on a
    test corpus to surface runtime failure modes that could not be anticipated
    from the paper: JSON parsing errors from LLM code-fence output, entity
    merge conflicts in Neo4j, and Windows encoding issues.
\end{enumerate}

\subsection{Modeling}

To make the requirements concrete before writing any code, the interactions
were mapped using a use case diagram (Figure~\ref{fig:usecase}). It answers
a straightforward question: who uses the system, and what do they actually
do with it? Three actors are involved --- the \emph{User}, the \emph{Ollama
LLM}, and \emph{Neo4j} --- and the diagram makes the boundary between
user-visible actions and internal system behaviour explicit at a glance.
Detailed functional requirements are listed in Section~3.3; data-flow
and structural models follow in Chapter~4.

% ── USE CASE DIAGRAM ──────────────────────────────────────────────
\begin{figure}[H]
\centering
\resizebox{0.96\textwidth}{!}{%
\begin{tikzpicture}[
  uc/.style={ellipse, draw=black, thin, fill=white,
             minimum width=3.8cm, minimum height=0.9cm,
             align=center, font=\small},
  conn/.style={-, thin},
  incl/.style={->, >=stealth, dashed, thin},
  lbl/.style={font=\scriptsize\itshape, fill=white, inner sep=1pt},
]
%% ── System boundary ─────────────────────────────────────────────
\draw[thick] (-0.4, 0.3) rectangle (11.9, 14.0);
\node[font=\small\bfseries] at (5.75, 13.75)
  {\textsc{TruthfulRAG v5}\ ---\ System Boundary};
\draw[thin,dashed,black!25] (5.75,13.4) -- (5.75,0.6);
\node[font=\tiny\itshape,black!50] at (2.8,13.3) {User-Initiated};
\node[font=\tiny\itshape,black!50] at (9.0,13.3) {System-Initiated};

%% ── Use cases: left column ──────────────────────────────────────
\node[uc] (sq)  at (2.8,12.4) {Submit Query};
\node[uc] (lc)  at (2.8,10.3) {Load Corpus};
\node[uc] (va)  at (2.8, 7.6) {View Answer};
\node[uc] (bat) at (2.8, 6.0) {Browse Audit Trail};

%% ── Use cases: right column ─────────────────────────────────────
\node[uc] (ga)  at (9.0,12.4) {Generate Final Answer};
\node[uc] (rrf) at (9.0,10.8) {Retrieve Ranked Facts};
\node[uc] (is)  at (9.0, 9.2) {Infer Schema};
\node[uc] (skg) at (9.0, 7.6) {Store Knowledge Graph};

%% ── Stick figure: User ──────────────────────────────────────────
\draw[thin] (-2.8,10.05) circle(0.27);
\draw[thin] (-2.8,9.78)  -- (-2.8,9.32);
\draw[thin] (-2.8,9.62)  -- (-3.05,9.42);
\draw[thin] (-2.8,9.62)  -- (-2.55,9.42);
\draw[thin] (-2.8,9.32)  -- (-2.98,9.00);
\draw[thin] (-2.8,9.32)  -- (-2.62,9.00);
\node[font=\small\bfseries] at (-2.8,8.72) {User};

%% ── Stick figure: Ollama LLM ────────────────────────────────────
\draw[thin] (13.6,11.35) circle(0.27);
\draw[thin] (13.6,11.08) -- (13.6,10.62);
\draw[thin] (13.6,10.92) -- (13.35,10.72);
\draw[thin] (13.6,10.92) -- (13.85,10.72);
\draw[thin] (13.6,10.62) -- (13.42,10.30);
\draw[thin] (13.6,10.62) -- (13.78,10.30);
\node[font=\small, align=center] at (13.6,10.02)
  {\textbf{Ollama}\\LLM};

%% ── Stick figure: Neo4j ─────────────────────────────────────────
\draw[thin] (13.6, 8.55) circle(0.27);
\draw[thin] (13.6, 8.28) -- (13.6, 7.82);
\draw[thin] (13.6, 8.12) -- (13.35,7.92);
\draw[thin] (13.6, 8.12) -- (13.85,7.92);
\draw[thin] (13.6, 7.82) -- (13.42,7.50);
\draw[thin] (13.6, 7.82) -- (13.78,7.50);
\node[font=\small, align=center] at (13.6,7.22)
  {\textbf{Neo4j}\\Graph DB};

%% ── Actor → use-case connections ────────────────────────────────
\draw[conn] (-2.53,10.05) -- (sq.west);
\draw[conn] (-2.53,10.05) -- (lc.west);
\draw[conn] (-2.53,10.05) -- (va.west);
\draw[conn] (13.33,11.35) -- (ga.east);
\draw[conn] (13.33,11.35) -- (is.east);
\draw[conn] (13.33, 8.55) -- (rrf.east);
\draw[conn] (13.33, 8.55) -- (skg.east);

%% ── <<include>> / <<extend>> ────────────────────────────────────
\draw[incl] (sq.east)  to node[lbl,above]{$\ll$include$\gg$} (ga.west);
\draw[incl] (sq.east)  to node[lbl,below]{$\ll$include$\gg$} (rrf.north west);
\draw[incl] (lc.east)  to node[lbl,above]{$\ll$include$\gg$} (is.west);
\draw[incl] (lc.east)  to node[lbl,below]{$\ll$include$\gg$} (skg.north west);
\draw[incl] (va.south) to node[lbl,right]{$\ll$extend$\gg$}  (bat.north);

\end{tikzpicture}}
\caption{Use case diagram for TruthfulRAG v5. The primary actor (\emph{User})
initiates four interactions. \emph{Ollama LLM} (secondary actor) handles schema
inference and answer generation; \emph{Neo4j} handles graph storage and
retrieval. Dashed arrows show \textit{include} (mandatory sub-step) and
\textit{extend} (optional enrichment) relationships.}
\label{fig:usecase}
\end{figure}

% ----------------------------------------------------------------
\section{Software Requirement Specification (SRS)}

\subsection{Purpose of the SRS}

This SRS defines the complete behaviour of TruthfulRAG v5. It is the binding
specification against which the implementation in Chapter~5 and the test results
in Chapter~6 are evaluated.

\subsection{Functional Requirements}

\begin{enumerate}
    \item[\textbf{FR1}] The system shall load documents and queries from a JSON
    corpus file. The JSON schema requires a \texttt{"docs"} list and optionally
    a \texttt{"queries"} list.

    \item[\textbf{FR2}] The system shall infer entity types and relation types
    from a configurable-size sample of the input documents using the LLM, without
    requiring any manual configuration.

    \item[\textbf{FR3}] The system shall extract subject-predicate-object triples
    from each document using the inferred schema, assigning an integer year to
    each triple where inferable.

    \item[\textbf{FR4}] The system shall maintain a count of how many distinct
    source documents assert each triple, storing this count as a property on the
    corresponding graph edge.

    \item[\textbf{FR5}] The system shall merge entity nodes whose names are at
    least 85\% similar (Ratcliff/Obershelp metric), using the Neo4j APOC library.

    \item[\textbf{FR6}] The system shall normalise all relation labels to
    \texttt{UPPER\_SNAKE\_CASE} via an LLM-assisted canonicalisation pass.

    \item[\textbf{FR7}] The system shall apply exponential temporal decay
    ($ e^{-\lambda t} $, $\lambda=0.08$) to edge scores during retrieval.

    \item[\textbf{FR8}] The system shall fuse BM25 keyword ranking and semantic
    embedding ranking using Reciprocal Rank Fusion when selecting which documents
    to focus on.

    \item[\textbf{FR9}] The system shall detect an explicit year in a user query
    via regular expression and, if found, exclude all graph edges with a year
    greater than the detected year.

    \item[\textbf{FR10}] The system shall detect contradicting paths (same
    relation-tail, different head entities), retain the path with the greater
    year, and discard the other, recording the reason.

    \item[\textbf{FR11}] When contradictions were removed, the system shall
    produce an explanation chain listing selected and rejected facts with reasons,
    and use this chain in the final answer prompt.

    \item[\textbf{FR12}] The system shall compute a 0--100\% confidence score
    per answer, combining entropy resolution signal, source support, and
    information recency according to configurable weights.
\end{enumerate}

\subsection{Non-Functional Requirements}

\begin{enumerate}
    \item[\textbf{NFR1}] \textbf{Privacy}: All computation shall run locally.
    No user documents shall be transmitted to external servers.

    \item[\textbf{NFR2}] \textbf{Domain Agnosticism}: The system shall function
    correctly on any text domain without code modification.

    \item[\textbf{NFR3}] \textbf{Robustness}: All LLM JSON outputs shall be
    parsed after stripping markdown code fences. Invalid JSON shall be handled
    gracefully without raising an unhandled exception.

    \item[\textbf{NFR4}] \textbf{Configurability}: All numeric hyperparameters
    shall be overridable via environment variables or the central \texttt{CFG}
    dictionary.

    \item[\textbf{NFR5}] \textbf{Efficiency}: Repeated identical LLM calls shall
    be served from an in-memory cache indexed by MD5 hash of (prompt, temperature).
\end{enumerate}

\subsection{Assumptions and Dependencies}

Three technical assumptions underlie the deployment model:

\begin{enumerate}
    \item Neo4j Community Edition is running locally with the APOC plugin
    installed. Without APOC, entity disambiguation and relation normalisation
    are silently skipped; the graph still functions but without those quality
    passes.
    \item Ollama is running locally with the configured model already pulled.
    The system does not attempt to download models at runtime.
    \item The Python environment contains all packages listed in
    Table~\ref{tab:deps}. No internet access is required after initial setup.
\end{enumerate}

\subsection{Constraints}

Three hard constraints shape the current implementation:

\begin{enumerate}
    \item Graph queries are limited to 200 nodes and 200 edges per retrieval
    call. For the six-document corpora tested here the cap is never reached,
    but it would need revisiting for large production corpora.
    \item Entropy sampling is fixed at $n=3$ samples to keep latency
    manageable on CPU-only hardware. Higher $n$ improves the entropy estimate
    but multiplies inference time accordingly.
    \item Entity disambiguation uses string similarity only
    (Ratcliff/Obershelp, 85\% threshold). Semantic entity linking against a
    knowledge base like Wikidata is not implemented in this version.
\end{enumerate}

% ----------------------------------------------------------------
\section{Software Requirement Validation}

\subsection{Review Process}

Each functional requirement was traced to a specific method in
\texttt{enhanced\_main.py} and verified by running the pipeline and
observing the relevant terminal output.

\begin{table}[H]
\centering
\caption{Requirement-to-implementation traceability}
\label{tab:req_trace}
\resizebox{\textwidth}{!}{%
\small
\begin{tabular}{|l|l|l|}
\hline
\textbf{Req.} & \textbf{Method in enhanced\_main.py} & \textbf{Verified By} \\
\hline
FR1  & \texttt{\_load\_corpus()}          & ``Loaded N docs'' line in terminal \\
FR2  & \texttt{\_infer\_schema()}         & Entity/relation list printed at start \\
FR3  & \texttt{\_extract()}               & Triples visible in Neo4j browser \\
FR4  & \texttt{\_accumulate()}            & Support count shown on graph edges \\
FR5  & \texttt{\_disambiguate()}          & ``Disambiguated N pairs'' in output \\
FR6  & \texttt{\_normalize()}             & ``Normalised N relation types'' line \\
FR7  & \texttt{\_temporal\_decay()}       & Older paths score lower in results \\
FR8  & \texttt{HybridRetriever}           & ``BM25+Semantic RRF'' in feature table \\
FR9  & \texttt{\_extract\_query\_year}   & Snapshot year logged when detected \\
FR10 & \texttt{\_detect\_contradictions}  & ``Contradictions removed: N'' line \\
FR11 & \texttt{\_build\_chain()}          & Explanation text appears in answer \\
FR12 & \texttt{\_confidence()}            & ``Confidence: XX\%'' printed at end \\
\hline
\end{tabular}}
\end{table}

\subsection{Validation Approach}

Validation was done by inspection of terminal output during live pipeline runs
and by checking that no unhandled exceptions were raised. All 12 functional
requirements produced observable evidence of correct behaviour during the test
runs documented in Chapter~6.

\subsection{Completeness and Consistency}

All 12 functional requirements and 5 non-functional requirements have
corresponding implementation methods in the codebase. The corroboration weight,
temporal decay factor, and confidence weights are all combined in a single
path-scoring expression, so there is no double-counting between the different
scoring signals. The entropy-based path selection follows the algorithm
described in arXiv:2511.10375 \cite{truthfulrag2024}, with the six extensions
added on top.


